\documentclass[11pt,a4paper]{article}
\usepackage[utf8]{inputenc}
\usepackage[margin=0.85in]{geometry}
\usepackage{amsmath,amssymb,amsfonts}
\usepackage{graphicx}
\usepackage{booktabs}
\usepackage{xcolor}
\usepackage{fancyhdr}
\usepackage{titlesec}
\usepackage{float}
\usepackage{listings}
\usepackage{hyperref}

\definecolor{navyblue}{RGB}{0, 51, 102}
\definecolor{darkslate}{RGB}{47, 79, 79}
\definecolor{codebg}{RGB}{245, 245, 245}

\hypersetup{
    colorlinks=true,
    linkcolor=navyblue,
    citecolor=navyblue,
    urlcolor=navyblue,
    pdftitle={Baseline Probabilistic Models: Reversal Task Benchmarks},
    pdfauthor={Lead Computational Architect}
}

\lstset{
    backgroundcolor=\color{codebg},
    basicstyle=\ttfamily\footnotesize,
    breaklines=true,
    frame=single,
    framerule=0.5pt,
    rulecolor=\color{darkslate}
}

\pagestyle{fancy}
\fancyhf{}
\rhead{\textcolor{darkslate}{\small Baseline Probabilistic Cognitive Models}}
\lhead{\textcolor{darkslate}{\small Benchmark Report}}
\rfoot{\thepage}

\titleformat{\section}{\large\bfseries\color{navyblue}}{\thesection}{1em}{}
\titleformat{\subsection}{\bfseries\color{darkslate}}{\thesubsection}{1em}{}

\title{\Huge \textbf{\color{navyblue} Competing Probabilistic Cognitive Models}\\[0.4em]
\large Benchmark Analysis \& Mathematical Formulations for 2AFC Reversal Tasks}
\author{\textbf{Lead Computational Architect \& Antigravity Research Group}}
\date{August 15, 2026}

\begin{document}

\maketitle

\begin{abstract}
This report provides the formal mathematical definitions, R implementation code, parameter estimation results, and model comparison metrics for discrete-time cognitive baselines benchmarked against the \texttt{ExactRModel} continuous-time cerebellar reservoir. Evaluated on 16,029 empirical human trials across 128 subjects, we benchmark: 1) a Counterfactual Expected Value (EV) Rescorla-Wagner model, and 2) a Win-Stay / Lose-Shift (WSLS) Markov heuristic model. Additionally, we propose two advanced Bayesian architectures: a Hidden Markov Model (HMM) and the Hierarchical Gaussian Filter (HGF).
\end{abstract}

\vspace{0.8em}
\hrule
\vspace{0.8em}

\section{Model Formulations \& Mathematical Update Rules}

\subsection{Model A: Counterfactual Expected Value Rescorla-Wagner (EV-RW)}
Because reward magnitudes ($M_A, M_B$) are explicitly visible to participants, the agent tracks internal probability estimates $\hat{p}_t(a) \in [0, 1]$ rather than raw Q-values. 

When option $a_t \in \{1, 2\}$ is chosen and outcome $F_t \in \{0, 1\}$ is received:
\begin{enumerate}
    \item \textbf{Direct Prediction Error Update (Chosen Option $a_t$)}:
    \begin{equation}
    \delta_t = F_t - \hat{p}_t(a_t)
    \end{equation}
    \begin{equation}
    \hat{p}_{t+1}(a_t) = \hat{p}_t(a_t) + \alpha \cdot \delta_t
    \end{equation}
    where $\alpha \in [0, 1]$ is the direct learning rate.
    
    \item \textbf{Counterfactual Prediction Error Update (Unchosen Option $\neg a_t$)}:
    \begin{equation}
    \hat{p}_{t+1}(\neg a_t) = \hat{p}_t(\neg a_t) + \alpha_c \cdot \left[ (1 - F_t) - \hat{p}_t(\neg a_t) \right]
    \end{equation}
    where $\alpha_c \in [0, 1]$ is the counterfactual learning rate.
    
    \item \textbf{Expected Value Computation}:
    \begin{equation}
    EV_A(t) = \hat{p}_t(1) \times M_A(t), \quad EV_B(t) = \hat{p}_t(2) \times M_B(t)
    \end{equation}
    
    \item \textbf{Softmax Action Selection}:
    \begin{equation}
    P(\text{Choice}_t = 1) = \frac{1}{1 + \exp\left(-\beta_{\text{softmax}} \cdot [EV_A(t) - EV_B(t)]\right)}
    \end{equation}
\end{enumerate}

\subsection{Model B: Win-Stay / Lose-Shift (WSLS) Markov Heuristic}
The WSLS baseline relies on a 2-parameter heuristic transition rule:
\begin{equation}
P(\text{Choice}_t = \text{Choice}_{t-1} \mid F_{t-1} = 1) = P_{\text{win\_stay}}
\end{equation}
\begin{equation}
P(\text{Choice}_t \neq \text{Choice}_{t-1} \mid F_{t-1} = 0) = P_{\text{lose\_shift}}
\end{equation}

\section{Empirical Fitting Results \& Model Benchmarks}

Both models were fitted via Maximum Likelihood Estimation (\texttt{optim}) across all 128 human subjects.

\begin{table}[H]
\centering
\caption{Empirical Model Fitting \& Comparison Metrics (128 Human Participants, 16,029 Trials).}
\label{tab:model_comp}
\vspace{0.5em}
\small
\begin{tabular}{l c c c c c}
\toprule
\textbf{Model Architecture} & \textbf{Free Params} & \textbf{Mean NLL} & \textbf{Total AIC} & \textbf{Total BIC} & \textbf{Key Fitted Parameters} \\
\midrule
\textbf{Counterfactual EV-RW} & $3$ & $58.66$ & $15,783.8$ & $16,870.3$ & $\alpha=0.4769, \alpha_c=0.1775, \beta=1139.0$ \\
\textbf{Win-Stay Lose-Shift} & $2$ & $53.25$ & $14,144.9$ & $14,867.2$ & $P_{\text{ws}}=0.8736, P_{\text{ls}}=0.5406$ \\
\bottomrule
\end{tabular}
\end{table}

\section{Advanced Bayesian Expansion Proposals}

To capture reversal dynamics better than standard RL, we propose two advanced extensions:

\subsection{Expansion 1: Hidden Markov Model (HMM) Latent State Tracker}
Maintains a belief vector $\mathbf{b}_t = [P(S_t = 1), P(S_t = 2)]^T$ over latent reversal state $S_t \in \{1, 2\}$, updating beliefs via transition matrix $\mathbf{T} = \begin{bmatrix} 1-\gamma & \gamma \\ \gamma & 1-\gamma \end{bmatrix}$ where $\gamma \ll 0.1$ is the expected reversal probability.

\subsection{Expansion 2: Hierarchical Gaussian Filter (HGF)}
A 3-level Bayesian filter tracking environmental volatility $\vartheta_t$. Precision-weighted prediction errors dynamically adjust learning rates $\alpha(t)$, accelerating belief updates immediately following a reversal.

\end{document}
